%% kapa: exam=1 round=36 year=2025 period=1 form=A subject=civil_law subject_ko=민법 seq=26 number=26
%% source: Q-Net 감정평가사 기출문제 (cst00309) artlSeq=5231525 file=raw/1008/5231525/2025년 제36회 감정평가사 1차시험 1교시 문제지 원본.pdf
%% conversion: kordoc 4.24.51 -> tools/gen1.py (auto); review=proofread-sonnet; flags=none
\begin{question}{26}
甲으로부터 甲소유의 X토지를 임차한 乙이 임대차기간이 만료된 후에도 다른 협의 없이 X토지를 반환하지 않고 점유ㆍ사용하고 있다. 이에 관한 설명으로 옳지 않은 것은? (다툼이 있으면 판례에 따름)
\begin{choices}
  \item 乙의 X토지에 대한 점유는 타주점유이다.
  \item 甲은 X토지를 간접점유하고 있다.
  \item 乙은 甲으로부터 X토지를 인도받은 시점부터 현재까지 계속하여 점유한 것으로 추정된다.
  \item 만약 乙이 X토지에 자기 소유의 건물을 신축하였다면 乙이 甲에게 X토지에 대한 소유의 의사를 표시한 것으로 볼 수 있다.
  \item 만약 甲이 종중이고 대표권 없는 종원인 丙이 종중과 무관하게 사인의 자격에서 乙에게 X토지를 임대한 것이라면 甲은 X토지를 간접점유하고 있다고 볼 수 없다.
\end{choices}
\end{question}
\answer{4}
